\clearpage
\section{Copying coherence and complete successor determinism}
\label{sec:copying-successor-coherence}
Copying now completes successor coherence for the original combined machine.
For any two transitions $s\to t$ and $s\to u$, the entire successor states
$t$ and $u$ are recursively isomorphic, including their token positions.
The result covers independently chosen contraction-tree boundaries, all
copyable payload families, recursive renaming, plugging, and simultaneous
expansion. It lifts to the unbounded machine's representation relation.
No determinism or certificate premise is added to the transition relation.

The construction leaves Proposition~12.2's terms, the Figure~16 translation,
the original copyability condition, and every stability condition unchanged.
An isomorphic normal presentation of a payload is used only inside the proof.
Lambda bodies remain complete recursive graphs; their contents are neither
executed nor replaced by an interface summary when copying.

\termdef{copying-successor-coherence}
\emph{Copying successor coherence} means that all actual copying rewrites at
one source request produce recursively isomorphic complete successor states.
The key observation is that the literal reduct retains every original
vertex and edge. It changes the selected source attachment to a fresh weakening
wire and redirects the payload's ordered output attachments through the
original distributors. One fresh payload copy and the distributor machinery
are then added. Thus the result can be described relative to the entire source
graph, without making the chosen tree boundary part of the description.

The proof does not assert that two independently chosen contraction trees are
isomorphic or even have the same number of external inputs. Instead it proves
that each literal local reduct, and then each contextual reduct, realizes the
same source-relative incidence description, up to recursive payload isomorphism.
This distinguishes the result from the earlier transport theorem, which
assumed a supplied isomorphism of tree presentations.

The dependency chain is: payload normal forms and rigidity; unique active
source frame and payload endpoint; a presentation retaining the whole source;
a certificate for the literal reduct; transport through every constructor of
\texttt{Rewrite}; comparison of two independently derived certificates; and
complete successor determinism. Exact-cost consequences follow from the
already proved transport of executions and finality along recursive isomorphism.

Appendix~B.1 removes copying successor coherence and unbounded determinism up
to representation from its outstanding obligations.
Section~\ref{sec:copying-intrinsic-refocusing} subsequently proves refocusing
for copying and allocation. Refocusing for the remaining operations, specimen
locality, the focused-graph correspondence, template arguments, semantic
weakening and Figure~22 remain open. The estimated completion is roughly
55--70\% of mathematical effort, with low confidence. This range is a judgement
about the remaining proofs, not a fraction of declarations or completed list
entries; Proposition~12.2 is still unproved.

\Needspace{12\baselineskip}
\subsection{Normal presentations of copyable payloads}

\termdef{copying-payload-normal-form}
A \emph{copying payload normal form} is either a single primitive payload or
the literal lambda edge followed by its boxed body. The datatype of plain
codes distinguishes instance, constant, application, unary and binary
operations. It retains the precise constant or operation parameter.
The variable-interface sizes are respectively $0,0,2,1,2$; only an instance
has an atom output, with atom-interface size one. The primitive map selects
the original label, and the graph map constructs its original one-edge corolla
at depth $d$. These definitions do not restrict copying to passive operations.

\lean{Copying/CopyPayloadNormal}{8}{14}{CopyRule.PlainCode}

\lean{Copying/CopyPayloadNormal}{16}{19}{CopyRule.PlainCode.variables}

\lean{Copying/CopyPayloadNormal}{21}{23}{CopyRule.PlainCode.atoms}

\lean{Copying/CopyPayloadNormal}{25}{30}{CopyRule.PlainCode.primitive}

\lean{Copying/CopyPayloadNormal}{32}{38}{CopyRule.PlainCode.net}

The primitive map is injective: case analysis on two codes reduces label
equality to equality of their constructors and parameters. No plain code has
the lambda primitive. These facts will exclude a plain/lambda ambiguity at a
shared payload head.

\lean{Copying/CopyPayloadNormal}{40}{41}{CopyRule.PlainCode.primitive_injective}

\lean{Copying/CopyPayloadNormal}{42}{42}{CopyRule.PlainCode.not_lambda}

The distinguished plain edge is the sole edge of its corolla. Its label and
source incidence prove that it is the head at input zero. Its source list is
exactly that singleton; its target list enumerates all environment outputs in
order. The proofs inspect the depth representation and primitive constructor;
none rearranges the boundary.

\lean{Copying/CopyPayloadNormal}{44}{44}{CopyRule.PlainCode.edge}

\lean{Copying/CopyPayloadNormal}{45}{46}{CopyRule.PlainCode.head}

\lean{Copying/CopyPayloadNormal}{47}{47}{CopyRule.PlainCode.source}

\lean{Copying/CopyPayloadNormal}{48}{48}{CopyRule.PlainCode.target}

The indexed family of normal forms records the graph just described, including
the full body parameter in its lambda constructor. Every graph satisfying the
existing \texttt{Copyable} predicate is recursively isomorphic to one of these
forms. Induction on copyability is direct in its six generating cases; in a
renaming case, compose the inverse supplied renaming with the induction
hypothesis. The existence of a normal form is therefore derived from the
original rule premise.

\lean{Copying/CopyPayloadNormal}{50}{52}{CopyRule.NormalForm}

\lean{Copying/CopyPayloadNormal}{54}{66}{CopyRule.normal_form}

For a normal form, the head is its primitive edge. The tail is that same edge
in a plain payload and the box edge in a lambda payload. The recorded primitive
is the plain label or lambda. The head is incident to input zero, its entire
source list is that singleton, and the tail targets enumerate the environment
outputs. For lambda, the latter equation follows through the right serial
insertion of the box. Every normal form satisfies the original copyability
predicate, by its corresponding constructor.

\lean{Copying/CopyPayloadNormal}{67}{69}{CopyRule.NormalForm.head}

\lean{Copying/CopyPayloadNormal}{70}{72}{CopyRule.NormalForm.tail}

\lean{Copying/CopyPayloadNormal}{73}{75}{CopyRule.NormalForm.primitive}

\lean{Copying/CopyPayloadNormal}{77}{80}{CopyRule.NormalForm.head_at}

\lean{Copying/CopyPayloadNormal}{81}{84}{CopyRule.NormalForm.head_source}

\lean{Copying/CopyPayloadNormal}{85}{92}{CopyRule.NormalForm.tail_target}

\lean{Copying/CopyPayloadNormal}{94}{102}{CopyRule.NormalForm.copyable}

\Needspace{12\baselineskip}
\subsection{Rigidity of two payload occurrences}

Occurrence embeddings compose by composing their port, internal-vertex and
edge maps. Injection and label preservation compose, and the ordered incidence
equations follow from the composition law for list mapping. The right component
of a serial graph has its canonical occurrence embedding: internal vertices
and edges receive the right-component tag, while boundary ports are attached
by serial insertion. These maps are injective on internal vertices and edges;
boundary ports are allowed to be attached by the context.

\lean{Copying/CopyPayloadRigidity}{8}{19}{Opening.Occurrence.comp}

\lean{Copying/CopyPayloadRigidity}{21}{35}{Opening.Occurrence.serialRight}

The relation \texttt{PayloadRelated} simultaneously asserts equality of the
variable and atom interface counts and recursive isomorphism of the complete
payloads. Its constructor uses the same counts on both sides, so elimination
also identifies the dependent interface types.

\lean{Copying/CopyPayloadRigidity}{38}{40}{CopyRule.PayloadRelated}

A literal lambda payload has one internal thunk wire between its head and box
tail. The head's target and tail's source are both the singleton containing
this wire. The definition and equations follow directly from serial composition.

\lean{Copying/CopyPayloadRigidity}{44}{45}{CopyRule.NormalForm.lambdaWire}

\lean{Copying/CopyPayloadRigidity}{46}{47}{CopyRule.NormalForm.lambda_head_target}

\lean{Copying/CopyPayloadRigidity}{48}{49}{CopyRule.NormalForm.lambda_tail_source}

Suppose two normal payload occurrences in $G$ have the same head edge.
Their primitive labels agree. If both are plain, primitive injectivity identifies
the codes and the payloads; the tail equality is the given head equality.
A mixed plain/lambda case is impossible. If both are lambdas, the head's
singleton target identifies the image of the thunk wire. Uniqueness of the
edge consuming that wire identifies the box tails. Their recursive box-label
relation then supplies equal interface counts and an isomorphism of the entire
bodies. Serial composition with the common lambda corolla yields the payload
isomorphism. This proves both \texttt{PayloadRelated} and equality of the actual
tail-edge images, without assuming an agreement between the two body descriptions.

\lean{Copying/CopyPayloadRigidity}{51}{86}{CopyRule.NormalForm.related}

\Needspace{12\baselineskip}
\subsection{The uniquely determined source data}

The selected input can be the first or second source of its binary
contraction. The list \texttt{selectedSources} records this order with a
Boolean: false means $[v,u]$, and true means $[u,v]$.

\termdef{copying-successor-match}
A \emph{copying successor match} consists of a normal payload occurrence in
$G$, the active contraction edge, this source-position Boolean, the other
source port, and a proved copy-resolution trace from the selected port $v$ to
the payload head. The active edge has the term-contraction label and exactly
the recorded source list. The match records no particular contraction-tree
boundary. Its tail is the image of the normal payload's tail.

\lean{Copying/CopySuccessorMatch}{9}{10}{CopyRule.selectedSources}

\lean{Copying/CopySuccessorMatch}{12}{26}{CopyRule.SuccessorMatch}

\lean{Copying/CopySuccessorMatch}{28}{28}{CopyRule.SuccessorMatch.tail}

The active-edge label and source list establish its head-at-$v$ property.
Source-list duplicate freedom proves that $v$ and the other source differ.
For two matches, source uniqueness first identifies the active edge. Equality
of its two ordered source lists then identifies the Boolean and the other
source. Opposite Boolean cases would repeat $v$, contradicting duplicate
freedom.

\lean{Copying/CopySuccessorMatch}{30}{33}{CopyRule.SuccessorMatch.head_at}

\lean{Copying/CopySuccessorMatch}{35}{38}{CopyRule.SuccessorMatch.distinct}

\lean{Copying/CopySuccessorMatch}{40}{49}{CopyRule.SuccessorMatch.frame_unique}

The two copy-resolution traces have the same source port, so the previously
proved resolution uniqueness identifies their payload heads. Payload rigidity
then identifies interface counts, relates the complete payloads recursively,
and identifies their actual tails. The tail target list has exactly the
environment length, since occurrence embeddings preserve incidence lists and
the normal tail enumerates those outputs.

\lean{Copying/CopySuccessorMatch}{51}{55}{CopyRule.SuccessorMatch.payload_unique}

\lean{Copying/CopySuccessorMatch}{57}{58}{CopyRule.SuccessorMatch.output_length}

Any occurrence embedding of $G$ into a larger graph transports this match.
Compose the payload occurrence, map the active edge and source ports, preserve
the Boolean, and transport the resolution trace. Source-list equality follows
by mapping the recorded ordered pair; primitive preservation supplies the label.

\lean{Copying/CopySuccessorMatch}{60}{68}{CopyRule.SuccessorMatch.map}

\Needspace{12\baselineskip}
\subsection{A presentation retaining the complete source graph}

The output-index map is precisely the existing pair-boundary permutation,
with the existing interface equality cast. It sends an output of either copy
to the matching position in the two term blocks followed by the two atom
blocks. No new permutation convention is introduced.

\lean{Copying/CopySuccessorPresentation}{9}{12}{CopyRule.copyOutput}

The fresh internal vertices comprise a second payload's internal vertices,
one weakening wire, the ordered output-link wires of both copies, and the
internal vertices of the original environment distributor with multiplicity
two. The fresh edges comprise the second payload's edges, one weakening edge,
and that distributor's edges. These are disjoint sums, so each kind has a
separate tag.

\lean{Copying/CopySuccessorPresentation}{14}{15}{CopyRule.AddedVertex}

\lean{Copying/CopySuccessorPresentation}{16}{21}{CopyRule.AddedEdge}

The weakening and output-link ports name their respective fresh wires.
A fresh payload port maps internal vertices into the clone, its sole input to
the old selected source $v$, and its outputs to the left-copy links. A distributor
port maps internal vertices into its fresh component, inputs to the ordered
links, and outputs to the original tail's ordered target list. The proved
length equation justifies every target-list index.

\lean{Copying/CopySuccessorPresentation}{23}{23}{CopyRule.SuccessorMatch.weakPort}

\lean{Copying/CopySuccessorPresentation}{24}{25}{CopyRule.SuccessorMatch.linkPort}

\lean{Copying/CopySuccessorPresentation}{27}{30}{CopyRule.SuccessorMatch.copyPort}

\lean{Copying/CopySuccessorPresentation}{32}{35}{CopyRule.SuccessorMatch.distributorPort}

\termdef{copying-source-presentation}
The \emph{source-relative copying presentation} has ambient carriers
\[
 V_G+V_{\mathrm{added}},\qquad E_G+E_{\mathrm{added}}.
\]
Every retention predicate is true. Old labels are retained; fresh labels come
from the payload, weakening and original distributor. On old edges, sources
are unchanged except that the selected source of the active contraction is
replaced by the fresh weakening wire, in the same ordered position. Targets
are unchanged except that the tail's outputs now enter the right-copy links.
The fresh copy receives $v$, the weakening has empty source and its fresh wire
as target, and the distributor connects the links to the old output attachments.
All these equations concern entire ordered incidence lists.

\lean{Copying/CopySuccessorPresentation}{37}{67}{CopyRule.SuccessorMatch.presentation}

\termdef{copying-successor-certificate}
A \emph{complete copying certificate} records a source match, a realization
of its presentation by the actual successor graph, and the equation sending
the resumed token to the old image of $v$. A realization preserves labels
recursively, has exactly the indicated image, and preserves source and target
lists. The certificate is proof data to be constructed from every actual
rewrite; it is not an additional rule premise.

\lean{Copying/CopySuccessorPresentation}{70}{74}{CopyRule.SuccessorCertificate}

\Needspace{12\baselineskip}
\subsection{Carrier bijections for the literal reduct}

The vertex map from the literal reduct to the presentation sends the residual
tree and retained payload to their original source occurrences. It sends the
active copy to the fresh payload summand, the residual weakening wire to the
fresh weakening summand, the serial joining wires to the output-link summand,
and the distributor vertices to their own summand. The inverse reconstructs
the nested serial and parallel tags from each of these alternatives.

\lean{Copying/CopySuccessorCarriers}{9}{27}{CopyRule.successorVertices}

\lean{Copying/CopySuccessorCarriers}{29}{40}{CopyRule.successorVertexInverse}

The edge map follows the same assignment: retained tree and retained payload
edges go to the old graph, while active-copy, weakening and distributor edges
go to their fresh summands. Its inverse reconstructs the original nested edge
tags. The vertex and edge maps are bijections. Both inverse equations are
proved by cases on depth and the disjoint-sum/component tags; the output-link
case respects the existing interface casts.

\lean{Copying/CopySuccessorCarriers}{42}{57}{CopyRule.successorEdges}

\lean{Copying/CopySuccessorCarriers}{59}{68}{CopyRule.successorEdgeInverse}

\lean{Copying/CopySuccessorCarriers}{70}{101}{CopyRule.successorVertexBijection}

\lean{Copying/CopySuccessorCarriers}{103}{126}{CopyRule.successorEdgeBijection}

A two-element source list containing the selected port can be written as one
of the two ordered forms $[v,u]$ and $[u,v]$. Inspect its two entries and the
membership proof to select the Boolean and other port. The local match uses
this fact on the actual active contraction supplied by the tree. The payload
occurs as the right serial component, and the existing trace theorem proves
resolution through the tree to its head.

\lean{Copying/CopySuccessorLocalMatch}{8}{19}{CopyRule.selected_sources_exists}

\lean{Copying/CopySuccessorLocalMatch}{21}{43}{CopyRule.localSuccessorMatch}

Every target of a non-tail edge of a normal payload is internal. There is no
such edge in a plain payload; in a lambda payload the only non-tail edge is
the head and its only target is the thunk wire. This small fact is essential
when showing that the retained copy changes only the tail's external outputs.

\lean{Copying/CopySuccessorLocalMatch}{45}{61}{CopyRule.NormalForm.nonTail_targets_internal}

\Needspace{12\baselineskip}
\subsection{Ports and ordered incidences of the literal construction}

A left serial output and the matching right serial input become the same
internal joining wire, indexed by that interface position. These equations
are proved from the serial pattern at each depth. They are boundary equations,
not a graph-rewriting simplification.

\lean{Copying/CopySuccessorPorts}{8}{14}{CopyRule.serial_left_output}

\lean{Copying/CopySuccessorPorts}{16}{24}{CopyRule.serial_right_input}

The pair insertion maps a port into the original parallel pair with its
original boundary cast. The active insertion takes the left payload into the
outer serial composition. The retained insertion first enters the right
payload after the residual tree, then the right parallel component, and then
the outer composition. The tree insertion similarly passes through the residual
frame, retained serial component and outer pair. These are the literal nested
attachment maps of the existing copying rule.

\lean{Copying/CopySuccessorPorts}{26}{28}{CopyRule.pairInsert}

\lean{Copying/CopySuccessorPorts}{30}{31}{CopyRule.activeInsert}

\lean{Copying/CopySuccessorPorts}{32}{34}{CopyRule.retainedInsert}

\lean{Copying/CopySuccessorPorts}{35}{38}{CopyRule.treeInsert}

Pair outputs have the previously specified pair-boundary indices. The active
copy's input is external input zero. The retained residual input $i$ becomes
external input $i+1$, because input zero has been assigned to the active copy.
Unfolding the boundary permutation proves the equations, with finite-index
equality established by equality of their natural-number values.

\lean{Copying/CopySuccessorPorts}{40}{44}{CopyRule.pair_output}

\lean{Copying/CopySuccessorPorts}{46}{50}{CopyRule.pair_left_input}

\lean{Copying/CopySuccessorPorts}{52}{60}{CopyRule.pair_right_input}

After the vertex bijection, the active insertion agrees with the presentation's
fresh-copy port map. Internal ports are immediate; the input case uses the
left-input equation, and outputs use the pair-output and serial-joining laws.
Retained payload internal and input ports map to their original source
occurrences. Its outputs instead map to the right-copy links. These last
ports are precisely the attachments that will be reconnected by the distributor.

\lean{Copying/CopySuccessorPorts}{62}{75}{CopyRule.active_port}

\lean{Copying/CopySuccessorPorts}{77}{78}{CopyRule.retained_internal}

\lean{Copying/CopySuccessorPorts}{80}{84}{CopyRule.retained_input}

\lean{Copying/CopySuccessorPorts}{86}{91}{CopyRule.retained_output}

The residual tree's internal and output ports map to their old source
occurrences. Its selected input zero maps to the new weakening wire. The
residual-frame input $i+1$ maps to residual external input $i$, and the full
outer construction restores it to source external input $i+1$. Thus all
unselected tree inputs preserve their original order. The proofs unfold the
residual frame and use the pair-input equations.

\lean{Copying/CopySuccessorPorts}{93}{94}{CopyRule.tree_internal}

\lean{Copying/CopySuccessorPorts}{96}{100}{CopyRule.tree_output}

\lean{Copying/CopySuccessorPorts}{102}{103}{CopyRule.tree_zero}

\lean{Copying/CopySuccessorPorts}{105}{110}{CopyRule.residual_succ}

\lean{Copying/CopySuccessorPorts}{112}{122}{CopyRule.tree_succ}

The diversion map sends exactly $v$ to the weakening wire and every other
source port to its old occurrence. Mapping an old edge's source list by this
function gives the presentation source list. For the active edge this is its
ordered pair and distinctness theorem; for any other edge, incidence uniqueness
prevents it from containing $v$. The tree attachment map agrees with this
diversion: inspect internal, output, selected-input and other-input cases.

\lean{Copying/CopySuccessorSource}{8}{12}{CopyRule.divertedPort}

\lean{Copying/CopySuccessorSource}{14}{29}{CopyRule.SuccessorMatch.source_divert}

\lean{Copying/CopySuccessorSource}{31}{50}{CopyRule.tree_port}

The local tail targets are the original external outputs in order, since the
payload is the right serial component. Therefore the literal distributor
attachment, after the vertex bijection, agrees with the presentation's
attachment map: internal vertices and inputs are direct, and outputs use
this tail-target equation at the same finite index.

\lean{Copying/CopySuccessorSource}{52}{60}{CopyRule.local_tail_targets}

\lean{Copying/CopySuccessorSource}{62}{74}{CopyRule.distributor_port}

The retained tree-edge name follows its residual-frame, serial and parallel
tags. The added-weakening edge name follows the same enclosing tags but selects
the residual weakening instead.

\lean{Copying/CopySuccessorIncidences}{9}{10}{CopyRule.keptTreeEdge}

\lean{Copying/CopySuccessorIncidences}{11}{12}{CopyRule.addedWeakening}

The active copy's source and target lists are the original payload lists
mapped by active insertion. The retained copy has the analogous equations
with retained insertion. The retained tree has them with tree insertion.
Each proof composes the literal list maps in the construction, so ordering
and repeated use of boundary attachments are preserved exactly. The new
weakening has empty source and the singleton residual image of the selected
tree input as target.

\lean{Copying/CopySuccessorIncidences}{14}{18}{CopyRule.active_source}

\lean{Copying/CopySuccessorIncidences}{20}{24}{CopyRule.active_target}

\lean{Copying/CopySuccessorIncidences}{26}{30}{CopyRule.retained_source}

\lean{Copying/CopySuccessorIncidences}{32}{36}{CopyRule.retained_target}

\lean{Copying/CopySuccessorIncidences}{38}{42}{CopyRule.kept_tree_source}

\lean{Copying/CopySuccessorIncidences}{44}{48}{CopyRule.kept_tree_target}

\lean{Copying/CopySuccessorIncidences}{50}{50}{CopyRule.added_weakening_source}

\lean{Copying/CopySuccessorIncidences}{51}{52}{CopyRule.added_weakening_target}

\Needspace{12\baselineskip}
\subsection{The complete local certificate}

The local match's tail is the right serial image of the payload tail. No
payload edge is the active contraction: the established outer-label
classification of copyable graphs excludes term-contraction labels there.
These facts separate the old-tree and old-payload cases of the incidence proof.

\lean{Copying/CopySuccessorLocal}{10}{10}{CopyRule.local_tail_eq}

\lean{Copying/CopySuccessorLocal}{12}{17}{CopyRule.local_active_not_payload}

For every ambient edge, take its inverse image in the actual reduct and map
its entire source list by the vertex bijection. The result is exactly the
presentation source list. Old tree edges use diversion; old payload edges
use retained insertion and cannot be active; new payload edges use active
insertion; weakening and distributor edges use their attachment equations.
The target proof has the same cases, but distinguishes the payload tail from
other payload edges. Tail targets become right-copy links, and non-tail
targets remain internal. This proves both incidence laws without discarding
any residual tree or host occurrence.

\lean{Copying/CopySuccessorLocal}{19}{57}{CopyRule.local_source}

\lean{Copying/CopySuccessorLocal}{59}{110}{CopyRule.local_target}

The local realization combines the two carrier bijections, exact image
properties, wire labels, recursively related edge labels, and the two ordered
incidence theorems. Exact image is surjectivity because all retention predicates
are true. The local certificate uses this realization and the local source
match. Its token equation holds because input zero maps to its old external
port. The entire original \texttt{after} graph is thus certified.

\lean{Copying/CopySuccessorLocal}{112}{156}{CopyRule.localRealization}

\lean{Copying/CopySuccessorLocal}{158}{162}{CopyRule.local_certificate}

\Needspace{12\baselineskip}
\subsection{Transport through occurrences and recursive renaming}

For an occurrence $f:G\hookrightarrow N$, the origin-port map sends old
ports through $f$ and fixes each kind of newly added port. The origin-edge map
similarly sends old edges through $f$ and fixes fresh edges. On old ports,
this agrees with ordinary occurrence attachment followed by old-source
inclusion. These definitions allow the source presentation to move with a
context without changing its fresh component.

\lean{Copying/CopySuccessorOccurrence}{10}{15}{CopyRule.SuccessorMatch.originPort}

\lean{Copying/CopySuccessorOccurrence}{17}{19}{CopyRule.SuccessorMatch.originEdge}

\lean{Copying/CopySuccessorOccurrence}{21}{23}{CopyRule.SuccessorMatch.origin_old}

Fresh-copy attachments commute with the origin-port map by cases on the
payload port. Distributor attachments commute too: their output case uses the
occurrence's target-list equation, list indexing through a map, and the proved
output length. This argument permits boundary attachment by the host.

\lean{Copying/CopySuccessorOccurrence}{25}{30}{CopyRule.SuccessorMatch.copyPort_map}

\lean{Copying/CopySuccessorOccurrence}{32}{43}{CopyRule.SuccessorMatch.distributorPort_map}

Every source and target list of the presentation commutes with these origin
maps. For old edges, injectivity of the occurrence edge map preserves the
active-edge and tail-edge tests. For fresh payload and distributor edges use
the attachment equations. The weakening case is immediate. Thus the complete
incidence description is natural under occurrences, not just under bijections.

\lean{Copying/CopySuccessorOccurrence}{45}{64}{CopyRule.SuccessorMatch.source_map}

\lean{Copying/CopySuccessorOccurrence}{66}{85}{CopyRule.SuccessorMatch.target_map}

For recursive isomorphism, the origin-port map is the port extension of the
old-vertex bijection plus identity on fresh vertices. This gives an isomorphism
of presentations, with the corresponding sum bijection on edges. Wire labels,
recursive edge labels and ordered incidences are preserved by the preceding
laws. A certificate can therefore be renamed on both source and result:
transport its presentation, rename its realization, and compose the port
equations for its token.

\lean{Copying/CopySuccessorOccurrence}{87}{92}{CopyRule.SuccessorMatch.origin_iso}

\lean{Copying/CopySuccessorOccurrence}{94}{117}{CopyRule.SuccessorMatch.renamePresentation}

\lean{Copying/CopySuccessorOccurrence}{119}{130}{CopyRule.SuccessorCertificate.rename}

\Needspace{12\baselineskip}
\subsection{Transport through a single plugged context}

The plugged vertex map retains host vertices as old host vertices; on an
inserted successor it uses the local realization and shifts its old tags into
the plugged source graph, leaving the fresh tags fresh. The edge map retains
all non-hole host edges and treats inserted edges analogously. Disjoint host,
old-component and fresh-component tags, together with local injection, prove
both maps injective.

\lean{Copying/CopySuccessorPlug}{13}{15}{CopyRule.SuccessorCertificate.plugVertices}

\lean{Copying/CopySuccessorPlug}{17}{19}{CopyRule.SuccessorCertificate.plugEdges}

\lean{Copying/CopySuccessorPlug}{21}{26}{CopyRule.SuccessorCertificate.plug_vertices_injective}

\lean{Copying/CopySuccessorPlug}{28}{33}{CopyRule.SuccessorCertificate.plug_edges_injective}

Both maps have the entire ambient presentation as their image. For a host
occurrence choose that occurrence itself; for an old or fresh component
occurrence choose a preimage under the local realization. The omitted hole
is absent on both sides of plugging. Hence the image equations agree with
the presentation's universally true retention predicates.

\lean{Copying/CopySuccessorPlug}{35}{42}{CopyRule.SuccessorCertificate.plug_vertex_range}

\lean{Copying/CopySuccessorPlug}{44}{51}{CopyRule.SuccessorCertificate.plug_edge_range}

Host ports commute with old-source inclusion by cases on their port kind.
Inserted component ports commute with the local realization followed by the
origin-port map; boundary cases reduce to host attachments, and internal
cases follow the disjoint tags. The source and target equations for inserted
ambient edges then follow from presentation naturality under the plugging
occurrence.

\lean{Copying/CopySuccessorPlug}{53}{55}{CopyRule.SuccessorCertificate.plug_host_port}

\lean{Copying/CopySuccessorPlug}{57}{70}{CopyRule.SuccessorCertificate.plug_component_port}

\lean{Copying/CopySuccessorPlug}{72}{76}{CopyRule.SuccessorCertificate.plug_source}

\lean{Copying/CopySuccessorPlug}{78}{82}{CopyRule.SuccessorCertificate.plug_target}

The plugged realization uses these maps, injection, image and port laws.
For a surviving host edge its original incidence lists are untouched; it cannot
be the active edge or tail because those belong to the inserted component.
For an inserted edge, use local realization followed by occurrence naturality.
Labels are preserved in the same two cases. The resulting plugged certificate
uses the transported source match and maps its resumed token through the host
attachment. This proves closure for the actual plugging constructor.

\lean{Copying/CopySuccessorPlug}{85}{131}{CopyRule.SuccessorCertificate.plugRealization}

\lean{Copying/CopySuccessorPlug}{133}{140}{CopyRule.SuccessorCertificate.plug}

\Needspace{12\baselineskip}
\subsection{Comparison of independently described payloads}

A recursive payload isomorphism acts on fresh payload vertices and edges;
it fixes the weakening, link and distributor summands. Its port extension
commutes with fresh-copy attachments, and with distributor attachments when
the two actual tail edges agree. The latter hypothesis ensures that both
presentations reconnect to the same ordered host targets.

\lean{Copying/CopySuccessorComparison}{10}{10}{CopyRule.SuccessorMatch.addedVertices}

\lean{Copying/CopySuccessorComparison}{11}{11}{CopyRule.SuccessorMatch.addedEdges}

\lean{Copying/CopySuccessorComparison}{13}{19}{CopyRule.SuccessorMatch.copyPort_congr}

\lean{Copying/CopySuccessorComparison}{21}{29}{CopyRule.SuccessorMatch.distributorPort_congr}

The presentation comparison is identity on the complete old graph and the
just-defined bijection on the fresh part. Match-frame uniqueness identifies
the active edge, source position and other port. Tail equality identifies
the redirected old target list. The payload isomorphism checks cloned labels
and incidence lists recursively; the distributor and weakening components
are unchanged. This constructs an isomorphism of the entire presentations.

\lean{Copying/CopySuccessorComparison}{31}{86}{CopyRule.SuccessorMatch.payloadPresentation}

For two certificates at one source request, payload uniqueness supplies
matching interface counts, the recursive payload isomorphism and tail equality.
Transport the first realization along the presentation comparison. It and the
second realization now have exactly the same ambient image, so the general
realization-comparison theorem constructs their recursive graph isomorphism.
Both token equations name the old selected port, which the comparison fixes.
The result is a state isomorphism, not merely agreement of payload heads.

\lean{Copying/CopySuccessorComparison}{88}{99}{CopyRule.SuccessorCertificate.coherent}

\Needspace{12\baselineskip}
\subsection{Transport through simultaneous expansion}

The existing dependent-fiber map uses a selected map at the changed filling
and bijections at every other filling. If the selected map is surjective,
so is the whole map into old fibers plus fresh occurrences. An old selected
fiber uses the selected preimage; another old fiber uses its bijection inverse;
a fresh occurrence also uses the selected preimage. This general lemma is
used for both vertices and edges.

\lean{Copying/CopySuccessorExpansion}{7}{25}{Opening.Fiber.surjective}

For an unchanged filling, equality of source and target fillings supplies
a recursive identity-transport isomorphism in the required direction.
The vertex and edge fiber maps use the certificate at the selected hole and
these isomorphisms elsewhere. The full expansion vertex map additionally
retains host vertices; the expansion edge map is the edge fiber map, since
expansion replaces every host edge by its filling.

\lean{Copying/CopySuccessorExpansion}{27}{28}{CopyRule.SuccessorCertificate.otherIso}

\lean{Copying/CopySuccessorExpansion}{30}{31}{CopyRule.SuccessorCertificate.fiberVertices}

\lean{Copying/CopySuccessorExpansion}{32}{33}{CopyRule.SuccessorCertificate.fiberEdges}

\lean{Copying/CopySuccessorExpansion}{34}{36}{CopyRule.SuccessorCertificate.expansionVertices}

\lean{Copying/CopySuccessorExpansion}{37}{37}{CopyRule.SuccessorCertificate.expansionEdges}

Fiber injection proves expansion-edge injection. Expansion-vertex injection
also distinguishes the retained host summand from all fiber and fresh summands.
Fiber surjectivity and local exact images prove the corresponding whole-image
properties; host vertices again supply their own preimages. Every ambient
occurrence is therefore represented exactly once.

\lean{Copying/CopySuccessorExpansion}{39}{44}{CopyRule.SuccessorCertificate.expansion_vertices_injective}

\lean{Copying/CopySuccessorExpansion}{45}{47}{CopyRule.SuccessorCertificate.expansion_edges_injective}

\lean{Copying/CopySuccessorExpansion}{49}{62}{CopyRule.SuccessorCertificate.expansion_vertex_range}

\lean{Copying/CopySuccessorExpansion}{64}{71}{CopyRule.SuccessorCertificate.expansion_edge_range}

There are three port equations. Host ports commute with old inclusion.
Ports in the selected successor filling commute with its realization and
source occurrence. Ports in any other filling commute with that filling's
unchanged isomorphism and then its old inclusion. Internal cases follow the
fiber maps; input and output cases reduce to the common host attachment.

\lean{Copying/CopySuccessorExpansion}{73}{75}{CopyRule.SuccessorCertificate.expansion_host_port}

\lean{Copying/CopySuccessorExpansion}{77}{87}{CopyRule.SuccessorCertificate.expansion_selected_port}

\lean{Copying/CopySuccessorExpansion}{89}{97}{CopyRule.SuccessorCertificate.expansion_other_port}

Source and target lists of the selected ambient component commute with its
fiber tag by occurrence naturality. For a different fiber, sources and targets
remain the original filling lists attached to the host. The active and tail
edge tests cannot select those edges: equality would identify two distinct
fiber tags. These four equations separate changed and unchanged fillings in
the whole incidence proof.

\lean{Copying/CopySuccessorExpansion}{99}{104}{CopyRule.SuccessorCertificate.expansion_source}

\lean{Copying/CopySuccessorExpansion}{106}{111}{CopyRule.SuccessorCertificate.expansion_target}

\lean{Copying/CopySuccessorExpansion}{113}{119}{CopyRule.SuccessorCertificate.outside_source}

\lean{Copying/CopySuccessorExpansion}{121}{127}{CopyRule.SuccessorCertificate.outside_target}

The expansion realization combines these maps and image laws. At the selected
fiber, labels and incidences follow the local realization and selected-port
law; at another fiber, they follow its equality-transport isomorphism and
other-port law. The expansion certificate adds the transported source match
and token equation. No assumptions about the other fillings' computations
are introduced; their equality is exactly the existing expansion constructor's
premise.

\lean{Copying/CopySuccessorExpansion}{129}{193}{CopyRule.SuccessorCertificate.expansionRealization}

\lean{Copying/CopySuccessorExpansion}{195}{202}{CopyRule.SuccessorCertificate.expand}

\Needspace{12\baselineskip}
\subsection{From actual rewrites to complete machine determinism}

The proposition \texttt{SuccessorWitness} existentially hides the payload
counts and graph while retaining a complete certificate. Any two witnesses
at the same source request give isomorphic successor states by certificate
comparison. For an actual local copy rule, normalize its copyable payload,
construct the normal local certificate, and transport it back along the
proved isomorphisms of the literal before and after constructions. Thus an
arbitrary original copyable payload, not just its normal presentation,
obtains a witness.

\lean{Copying/CopyCoherence}{7}{10}{CopyRule.SuccessorWitness}

\lean{Copying/CopyCoherence}{12}{17}{CopyRule.SuccessorWitness.coherent}

\lean{Copying/CopyCoherence}{19}{28}{CopyRule.local_witness}

An actual rewrite already supplies a recognized head. If an occurrence
embedding exhibits a head at the mapped request, unique source incidence
identifies its primitive with that recognized head. This reflects existence
of the requested primitive head back to the source. For a local rewrite whose
head is a term contraction, all non-copy constructors contradict their own
proved operation heads. The remaining copy constructor supplies the local
witness just established.

\lean{Copying/CopyCoherence}{30}{36}{Combined.rewrite_head_exists_reflect}

\lean{Copying/CopyCoherence}{38}{50}{Combined.local_copy_certificate}

Induction on the original rewrite derivation now gives a copying witness
for every rewrite at a term-contraction head. The local case uses the local
certificate theorem; recursive renaming, plugging and expansion reflect the
head to the induction hypothesis and apply the corresponding certificate
transport. Consequently any two actual copying rewrites at that request have
recursively isomorphic whole results and corresponding search tokens.

\lean{Copying/CopyCoherence}{53}{68}{Combined.rewrite_copy_certificate}

\lean{Copying/CopyCoherence}{71}{74}{Combined.copying_rewrite_deterministic}

\termdef{combined-successor-determinism}
\emph{Complete successor determinism up to representation} asserts
\[
 s\to t\ \wedge\ s\to u\quad\Longrightarrow\quad
 \operatorname{Nonempty}(\operatorname{StateIso}(t,u)).
\]
The earlier exhaustive rewrite classification either already provides the
successor isomorphism or identifies a common copying head. The new copying
comparison resolves that final alternative, proving rewrite determinism.
For full steps, two search successors are literally equal, search and rewrite
cannot overlap in focus mode, and two rewrites use the completed theorem.
This is an unconditional theorem about the existing transition relation;
it makes no template, refocusing or termination assumption.

\lean{Copying/CopyCoherence}{77}{82}{Combined.rewrite_deterministic}

\lean{Copying/CopyCoherence}{85}{91}{Combined.step_deterministic}

For any actual step $s\to t$, successful termination in $n$ steps from $s$
is equivalent to successful termination in $m$ steps from $t$, with $n=m+1$.
The forward direction compares the first step of the successful run with the
chosen step, then transports its remaining run along the successor isomorphism.
The reverse direction prefixes the chosen step. Existing final-normality and
representation-transport lemmas justify the zero-length and finality cases.

\lean{Copying/CopyCoherence}{93}{95}{Combined.cost_spectrum}

Every unbounded step is a lift of one finite-depth step. Two such steps from
the same unbounded state have the same depth and interfaces, so the completed
finite theorem lifts to \texttt{Represents} between their successors. Previously
proved representation invariance then gives observational state equivalence
for every cost order $Q$, in particular exact costs. Literal equality of
named graph carriers is not claimed; recursive representation is the conclusion.

\lean{Copying/CopyCoherence}{97}{101}{Combined.unbounded_deterministic}

\lean{Copying/CopyCoherence}{103}{105}{Combined.unbounded_successors_equiv}

\Needspace{12\baselineskip}
\subsection{Exact costs of copying and of arbitrary finite prefixes}

\termdef{combined-exact-prefix-cost}
An \emph{exact prefix cost law} says that for any actual finite run
$s\steps{r}t$,
\[
 \operatorname{TerminatesIn}(n,s)
 \quad\Longleftrightarrow\quad
 \exists m\;(n=m+r\ \wedge\operatorname{TerminatesIn}(m,t)).
\]
Induct on the given prefix. The empty prefix is immediate; the successor
case uses the one-step spectrum theorem and adds its one transition to the
inductive cost. Apply this law to a successful run ending in a final state.
Final normality forces any successful continuation there to have length zero,
so every successful run from the same source has the same length. This
uniqueness concerns the entire combined machine, including copying.

\lean{Copying/CopyCost}{8}{20}{Combined.run_cost_spectrum}

\lean{Copying/CopyCost}{23}{32}{Combined.successful_cost_unique}

Every selected input of a genuine contraction tree requests copying in one
search transition: use its proved active term-contraction head and the existing
search-entry rule. The actual local copying step then gives the exact
one-transition cost equation from rewrite mode to the literal reduct in
search mode. Composing search entry with this equation gives two transitions
from the original search state to the reduct's continuation. These equations
also hold if no successful continuation exists.

\lean{Copying/CopyCost}{35}{39}{CopyRule.requests_copy}

\lean{Copying/CopyCost}{42}{46}{CopyRule.copying_cost_iff}

\lean{Copying/CopyCost}{49}{57}{CopyRule.initial_cost_iff}

Under an arbitrary serial host, transport this two-transition run using the
existing context law and apply the general prefix-cost theorem. The continuation
is the actual reduct attached at the same host position. In particular an
internal return is not silently declared final; its remaining host execution
is still counted.

\lean{Copying/CopyCost}{60}{67}{CopyRule.host_handoff}

The active fresh lambda returns in one search step without examining its
body. Its copied head label is lambda; its source list is the original head's
singleton mapped by active insertion, hence the selected input. The original
lambda entry rule returns there. This argument applies to arbitrary body
graphs and arbitrary nesting, independently of whether a later application
of the body would terminate.

\lean{Copying/CopyCost}{70}{82}{CopyRule.lambda_returns}

If the actual copied payload has this immediate return transition, then its
initial successful cost is exactly three: one search transition, one copying
transition, and the return. The exact two-step continuation law reduces the
claim to the return law and zero cost at the final external input. Constants
and instances use their existing return theorems; lambdas use the new theorem.
The resulting three cost equivalences quantify over every successful run of
the full machine and every matched contraction tree.

\lean{Copying/CopyCost}{85}{95}{CopyRule.passive_cost_iff}

\lean{Copying/CopyCost}{97}{100}{CopyRule.constant_cost_iff}

\lean{Copying/CopyCost}{102}{105}{CopyRule.instance_cost_iff}

\lean{Copying/CopyCost}{107}{110}{CopyRule.lambda_cost_iff}

\Needspace{12\baselineskip}
\subsection{Semantic regression scenarios and the remaining boundary}

The closed constant-copying program uses the existing one-input canonical
contraction tree and original constant corolla. Its focus- and hole-freedom
come from the established tree/payload label-safety theorem. Every successful
run of this program in the unbounded machine, at every finite padding, costs
exactly three. Exact padding reflection reduces the assertion to its original
depth, exact finite-run representation reduces it to the combined finite
machine, and the constant cost theorem finishes the proof. This graph is a
copying-rule regression; no claim identifies it with a new term translation.

\lean{Copying/CopyCoherenceExamples}{8}{14}{CopyRule.Examples.constantProgram}

\lean{Copying/CopyCoherenceExamples}{17}{21}{CopyRule.Examples.constant_pad_cost_iff}

Copying a lambda whose body is the unchanged translated nested-lambda body
also has exact cost three. This is an instance of the arbitrary-body theorem;
it ensures that copying and returning the outer function do not execute either
body. Separately, the two distinct copied reference instances in the existing
shared-store graph must resolve to the same cell, with three trace hops each.
Compare any proposed resolutions with the previously constructed traces by
store-resolution uniqueness. Thus the coherence argument remains compatible
with reference aliasing, rather than treating copying as store allocation.

\lean{Copying/CopyCoherenceExamples}{24}{27}{CopyRule.Examples.nested_lambda_cost_iff}

\lean{Copying/CopyCoherenceExamples}{30}{37}{CopyRule.Examples.shared_reference_unique}

These regressions prove copying costs, preservation of nested bodies and
store identity, and compatibility with the separately proved depth-padding
semantics. They do not constitute the Figure~22 contextual-equivalence proof,
a refocusing theorem, or robustness of the copying template. In particular,
repeated closure applications and divergent callbacks still have to be handled
in the general specimen and template arguments. The present milestone supplies
the unconditional successor-determinism law those arguments require.
